\section{Parity remains hard at nearly logarithmic learner depth}
\label{sec:critical-nearly-logarithmic-depth}

The comparison can be extended to a larger growing-depth range by
retaining the leading constant in both the interpolating competitor's
rate and the transport lower bound. The first improvement comes from
tracking the freely cooling feature variance during mean amplification.
It avoids charging a fixed loss at every hidden layer.

\begin{lemma}[Mean amplification with its leading constant]
\label{lem:critical-sharp-mean-amplification}
Fix $a\in(0,1]$. Suppose a first-feature law $\nu$ has entropy at most
$D\ge1$ relative to the unmodified first-feature law $\mu$, and its
mean $h$ satisfies $g(x)h(x)\ge a$ everywhere, for a Boolean $g$.
Put $q=(D/a^2)^{1/L}$. There are constants $K,P,C$, depending only on
$a$, such that, if $L\ge2$ and $q\ge KL^P$, then
\[
 J_L(g)\le\frac q2\{L+C\log L\}.
\]
The first law and the subsequent covariance diagonals need not be
constant over inputs.
\end{lemma}

\begin{proof}
Let $F(v)=\mathbb E\tanh^2(\sqrt vZ)$ and $u_s=F^s(2)$.
The function $F$ is nondecreasing and one-Lipschitz. Monotonicity is
pointwise in $v$, and Gaussian differentiation gives $F'\le1$ because
$(\tanh^2)''\le2$. On $[0,2]$, one has $v-F(v)\ge cv^2$ with
universal $c>0$: the ratio tends to two at zero and is strictly
positive away from zero. It follows that
\[
 u_1<1,\qquad u_s\le\frac{C_0}{s+1},\qquad
 \sum_{s=1}^L u_s\le C_1\log(L+1).
\]
For a centered Gaussian $G$ of variance $v$, the same second-derivative
bound and symmetry give
\[
 \mathbb E\tanh^2(G+t)\le F(v)+t^2.
\]
For $t\ge0$, integration of
$\mathbb E\operatorname{sech}^2(G+s)\ge1-v-s^2$ gives the convenient
weaker estimate
\[
 \mathbb E\tanh(G+t)\ge t(1-v-t^2).
\]

Set $\theta=q/(2D)=1/(2a^2q^{L-1})$.
Use $(1-\theta)\mu+\theta\nu$ at the first hidden layer. Its entropy
is at most $q/2$, and its second moment dominates $b_1b_1^T$ with
$b_1=\sqrt\theta h$. Its maximum diagonal is at most $u_1+\theta$.
For $L=2$, a readout proportional to $b_1$ costs at most $q/4$,
proving the assertion. Assume henceforth $L\ge3$.

Let $\kappa=1-u_1>0$, $\epsilon=\kappa/(8L^2)$, and
\[
 c_s=1-u_{s-1}-\epsilon\qquad(2\le s\le L-1).
\]
These constants are uniformly positive. The preceding bounds on $u_s$
imply, for universal constants $A_0,C_2$,
\[
 \prod_{s=2}^{L-1}c_s^{-1}\le L^{A_0},\qquad
 \sum_{s=2}^{L-1}c_s^{-2}\le L-2+C_2\log L.
\]
Indeed, $-\log c_s$ and $c_s^{-2}-1$ are both at most a universal
multiple of $u_{s-1}+\epsilon$.

At hidden layer $s$, shift the reference Gaussian by
\[
 m_s(x)=\frac{\sqrt q}{c_s}b_{s-1}(x),
\]
and let $b_s$ be its tanh feature mean. For $s=2$, $b_1$ denotes
the tracked component just constructed; thereafter $b_{s-1}$ is the
actual preceding mean. In either case its outer product is dominated
by the preceding second moment. Thus the shift lies in the covariance
range and costs at most $q/(2c_s^2)$, including singular references.

The centered-Gaussian mean $\mathbb E\tanh(G+t)$ is odd, has the sign
of $t$, and has absolute value at most $|t|$. Therefore, without yet
using a variance bound,
\[
 \max_x|m_s(x)|
 \le\frac{L^{A_0}}{\sqrt2a}\,q^{-(L-s)/2}.
\]
For $q\ge2$, summing the resulting geometric series gives
\[
 \theta+\sum_{s=2}^{L-1}\max_x|m_s(x)|^2
 \le C_a L^{2A_0}/q.
\]
Choose $K,P$ so that $q\ge KL^P$ makes this at most $\epsilon$;
the same choice ensures $\theta\le1$.
Since $F$ is one-Lipschitz, the shifted-square estimate gives inductively
\[
 \max_x Q_s(x,x)
 \le u_s+\theta+\sum_{t=2}^s\max_x|m_t(x)|^2.
\]
In particular the variance error before layer $s$, together with the
squared current shift, is at most $\epsilon$.
The shifted-mean estimate consequently yields
\[
 g(x)b_s(x)\ge c_s g(x)m_s(x)
             =\sqrt q\,g(x)b_{s-1}(x).
\]
The last hidden mean has global signed margin at least
\[
 \sqrt\theta a q^{(L-2)/2}=\frac1{\sqrt{2q}}.
\]
A readout proportional to this mean has output margin $1/2$ and costs
at most $q/4$. Adding all costs gives
\[
 \frac q2+\frac q2\sum_{s=2}^{L-1}c_s^{-2}+\frac q4
 \le\frac q2\{L+C\log L\},
\]
as required.
\end{proof}

\begin{lemma}[Joint optimization of the transport budget]
\label{lem:critical-joint-budget}
If a Boolean function depends on all $d$ original input coordinates,
then any admissible critical path of cost $A$ satisfies, for every $t>0$,
\[
 d\le8(2+1/t)^{L-2}
 \left(\frac{A+(L-2)t}{L-1}\right)^{L-1}
 \exp\!\left(2\sqrt{2A/d}\right).
\]
In particular, if $A\le(q/2)\{L+C\log L\}$, then
\[
 d\le4q^{L-1}
 \exp\!\left(C'\log L+C'L/\sqrt q+2\sqrt{2A/d}\right),
\]
where $C'$ depends only on $C$.
\end{lemma}

\begin{proof}
Use the sharper pathwise transport inequality
\eqref{eq:critical-transport-exact-budget}, with readout cost $r$,
hidden costs $D_s$, and $k\ge d$ relevant augmented coordinates.
Cauchy--Schwarz gives
\[
 (1+\sqrt{2D_s})^2\le(2+1/t)(D_s+t).
\]
The arithmetic-geometric mean inequality gives
\[
 r\prod_{s=2}^{L-1}(D_s+t)
 \le\left(\frac{A+(L-2)t}{L-1}\right)^{L-1}.
\]
The first-layer factor is at most $\exp(2\sqrt{2A/d})$.
Combining these facts proves the first display, also when $L=2$ and
the products are empty. For the second choose $t=\sqrt q$, factor out
$2^{L-2}$ and $(q/2)^{L-1}$, and use
$\log(1+u)\le u$. The resulting errors are bounded by
$C'\log L+C'L/\sqrt q$.
\end{proof}

\begin{theorem}[Critical parity lower bound at nearly logarithmic depth]
\label{thm:critical-nearly-logarithmic-depth}
Fix $0<\eta<1$. Let the target be full parity on $d$ coordinates and
let integer depths $L=L(d)\ge2$ satisfy
\[
 L=o\!\left(\frac{\log d}{\log\log d}\right).
\]
Put $M=d^{1/(L-1)}$, and draw
$1\le n\le M^{1-\eta}$ independent uniform cube inputs, labeled by
parity. For all sufficiently large $d$, with training probability at
least $1-1/(2M)$, the original globally margin-conditioned posterior at
$\gamma=1$ satisfies
\[
 \lim_{N\to\infty}\text{Gibbs error}=\frac12,
 \qquad
 \liminf_{N\to\infty}\Pr_x\{p_N(x)<2/3\}\ge\frac14.
\]
For each $d$, the width limit fixes $d,L,n$ first. Moreover, there is a
universal $c>0$ for which the same conclusion holds for every
\[
 2\le L\le c\eta\frac{\log d}{\log\log d},
\]
once $d$ is sufficiently large depending on $\eta$.
\end{theorem}

\begin{proof}
Let $m$ be the smallest odd integer at least
$12n\log(2M(n+1))$. For large $d$,
\[
 m\le50n\log M\le50M^{1-\eta}\log M=o(M)\le o(d).
\]
Lemma~\ref{lem:critical-oriented-majority} supplies, with the stated
training probability, a consistent Boolean function and a first-feature
law of entropy $D=dm$ and global signed mean at least $a=1/100$.
Put $q=(10000dm)^{1/L}$. The depth hypothesis implies
\[
 \log q\ge\frac{\log d}{L}\gg\log\log d,
\]
so $q$ exceeds every fixed power of $L$.
Lemma~\ref{lem:critical-sharp-mean-amplification} gives a consistent
competitor of rate at most $U=(q/2)\{L+C\log L\}$.

Suppose a Boolean function retaining all $d$ coordinates had rate at
most $U$. Apply Lemma~\ref{lem:critical-joint-budget} to paths whose
costs approach its infimum, then use continuity and monotonicity in
the budget. It follows that
\[
 \log\frac d{q^{L-1}}
 \le\log4+C'\log L+C'L/\sqrt q+2\sqrt{2U/d}.
\]
The last two terms tend to zero. For $L\ge3$, use $m\le d$,
$q\le(10000d^2)^{1/L}$, and $L=O(\log d)$ to see $U/d\to0$.
For $L=2$, use $m\le50d^{1-\eta}\log d$ instead.
On the other hand, the sample bound and $q^L=10000dm$ give
\[
 \log\frac d{q^{L-1}}
 \ge\frac{\eta\log d}{L}-\log(500000\log M).
\]
Here the first term dominates $\log\log d$, whereas the other
logarithmic terms are $O(\log\log d)$. This is a contradiction.
There is therefore a strict rate gap between the consistent competitor
and every Boolean function retaining all $d$ input coordinates.

The exact finite-domain posterior rate principle makes posterior mass
on the latter class tend to zero. Every remaining function omits a
parity coordinate and has uniform parity error exactly $1/2$.
This proves the Gibbs limit, and Markov's inequality gives the displayed
confidence-error bound from $\mathbb E_xp_N(x)\to1/2$.

For the final depth range, choose the universal $c$ small enough that
$\log q\ge\log d/L$ guarantees the fixed power of $L$ required by
the amplification lemma, and that
$\eta\log d/L$ exceeds $(C'+2)\log\log d$ by a fixed factor.
This is possible independently of $\eta\le1$.
The remaining estimates hold for all sufficiently large $d$ depending
on $\eta$, completing the proof.
\end{proof}

This theorem concerns the actual posterior, including mixtures of
minimizers. It does not reach the order-$\log^2d$ depths at which
controlled parity seeds attain a bounded absolute rate, and it does
not by itself give an exponential-in-teacher-depth lower bound.
